%   % !TEX root = ../../VIII,3_Rahmen-TeX_8-1.tex
%
%
%   Band VIII, 3 N.~??A33
%   Signatur/Tex-Datei: LH_35_10_06_037
%   RK-Nr. 60029 (mit RK 41772 und RK 58249 zusammen)
%   Überschrift: [Hypothesis de fluido elastico materiae interfuso ((ex: Quomodo agant corpora in se invicem))]
%   Modul: Mechanik / AEF (Elastizität)
%   Datierung: [nach April 1704]
%   WZ: Keins
%   SZ: Keins
%   Bilddateien (PDF): LH_35_10_06_037_d (insgesamt: eine)
%
%
\selectlanguage{ngerman}%
\frenchspacing%
%
\begin{ledgroupsized}[r]{120mm}%
\footnotesize%
\pstart%
\noindent\textbf{Überlieferung:}%
\pend%
\end{ledgroupsized}%
\begin{ledgroupsized}[r]{114mm}%
\footnotesize%
\pstart \parindent -6mm
\makebox[6mm][l]{\textit{L}}%
Aufzeichnung:
LH~XXXV~10,~6 Bl.~37.
Ein längs halbiertes Blatt 2\textsuperscript{o} (32 x 10,5 cm).
Eine Seite auf Bl.~37~r\textsuperscript{o}\!.
Die ersten einundzwanzig Textzeilen
(bis ungefähr S.~\refpassage{LH_35_10_06_037r_blstft_udhakv}{LH_35_10_06_037r_blstft_udhakv})
wurden mit Bleistift aufgesetzt und nachträglich mit Tinte überschrieben;
da die Bleistiftfassung weitgehend verblasst ist,
lässt sich deren Übereinstimmung mit der Tintenfassung nur man\-cher\-orts feststellen.
Bl.~37~v\textsuperscript{o}\! überliefert ein Fragment eines Brief\-kon\-zep\-tes von Leibnizens Hand,
welches voraussichtlich in \textit{LSB} II,~4 ediert wird.
Der Text\-trä\-ger von N.~33 hing ursprünglich mit der Handschrift LH~XXXV~14,~2 Bl.~60 zusammen,
die ein weiteres Fragment des letztgenannten Briefkonzeptes sowie eine voraussichtlich in \textit{LSB}~VII zu edierende Aufzeichnung über den Schwerpunkt überliefert.
\pend%
\end{ledgroupsized}%
%
\vspace{8mm}%
\begin{ledgroup}%
\footnotesize%
\pstart%
\noindent\footnotesize{\textbf{Datierungsgründe:}\label{LH_35_10_06_037_datierung_rhtj}
Das auf dem Träger des Textes N.~33 mit überlieferte Brieffragment gehört zum verworfenen Konzept eines Schreibens von Leibniz an W.~Pulteney (d.J.), das die Herausgeber von \textit{LSB} II,~4 auf die Zeit nach April 1704 datieren.\protect\index{Namensregister}{\textso{Pulteney}, William 1684\textendash1764}}
Das Brieffragment lag zweifelsohne bereits vor, als N.~33 verfasst wurde.
Hieraus ergibt sich der Terminus post quem der Datierung.
Zur Ermittlung eines Terminus ante quem gilt es die nur flüchtige Bekanntschaft zwischen Leibniz und Pulteney zu berücksichtigen.
Die Anzahl der Erwähnungen von Pulteney im Leibniz-Nachlass ist entsprechend begrenzt und reicht nach heutigem Kenntnisstand nicht über den Anfang 1705 hinaus (G.\,Guidis Brief an Leibniz vom 4. Januar, \textit{LSB} I,~24 N.~152, S.~259.13\cite{01342}).%
\protect\index{Namensregister}{\textso{Guidi} Giuseppe, gest. 1720 (?)}
Somit ist wahrscheinlich, dass auch das mit überlieferte Brieffragment nicht viel später entstand.
Da es nicht ausgeschlossen ist, dass Leibniz erst später die unbeschriebene Seite des Briefkonzeptes für die Aufzeichnung N.~33 verwendete, bleibt der vorgeschlagene Terminus ante quem aber hypothetisch.
\pend%
\end{ledgroup}%
%
\selectlanguage{latin}%
\frenchspacing%
%
%
\vspace{8mm}%
%
\count\Bfootins=1200%
\count\Afootins=1200%
\count\Cfootins=1200%
%
\pstart%
\noindent%
%
\lbrack37~r\textsuperscript{o}\rbrack\    %%%%    Blatt 37r
%
Si non contactu\protect\index{Sachverzeichnis}{contactus}
sed vi interpositi\protect\index{Sachverzeichnis}{vis elastri}
Elastri\protect\index{Sachverzeichnis}{elastrum interpositum}
corpus agat in Corpus\protect\index{Sachverzeichnis}{actio corporis in corpus} %
% \lbrack,\rbrack\
efficiemus phaenomena\protect\index{Sachverzeichnis}{phaenomenon}
non supponendo corpus quiescens\protect\index{Sachverzeichnis}{corpus quiescens}
quasi occurrere ictui alterius,\protect\index{Sachverzeichnis}{ictus corporis}
\edtext{sed}{%
\lemma{sed}\Bfootnote{%
\textit{erg.~L}}}
%
erunt resistentiae\protect\index{Sachverzeichnis}{resistentia} tamen
ut moles\protect\index{Sachverzeichnis}{moles}.
Si porro 
\edtext{concipiamus materiae\protect\index{Sachverzeichnis}{materia}}{%
\lemma{concipiamus}\Bfootnote{%
\textit{(1)}~materiam
\textit{(2)}~materiae%
~\textit{L}}}
%
fluidum elasticum\protect\index{Sachverzeichnis}{fluidum elasticum} aequaliter intelligi posse interfusum\lbrack,\rbrack\
licebit Massam\protect\index{Sachverzeichnis}{massa} intelligere ut
divisam in partes aequales,
nec aliud nisi vim elasticam\protect\index{Sachverzeichnis}{vis elastica}
utrinque in aequalia aequaliter agere
nullo discrimine quiescentis aut moti\lbrack,\rbrack\
obtinebitur etiam effectus aequivalens ei\protect\index{Sachverzeichnis}{effectus aequivalens}
qui evenisset sine mutatione\protect\index{Sachverzeichnis}{effectus sine mutatione}%
\lbrack,\rbrack\
cum scilicet aequale incurrit in aequale.
Addatur jam attractio corporum\protect\index{Sachverzeichnis}{attractio corporum} rursus
ex fluidi interpositi\protect\index{Sachverzeichnis}{fluidum interpositum} impulsu,\protect\index{Sachverzeichnis}{impulsus fluidi}
quae simul efficit,
ut cohaereant corpora,\protect\index{Sachverzeichnis}{cohaesio corporis}
et totum compositum motu centri gravitatis%
\protect\index{Sachverzeichnis}{motus centri gravitatis}\protect\index{Sachverzeichnis}{centrum gravitatis}
ex legibus aequilibrii\protect\index{Sachverzeichnis}{lex aequilibrii}\protect\index{Sachverzeichnis}{aequilibrium}
seu vis mortuae\protect\index{Sachverzeichnis}{vis mortua} feratur.
In his omnibus tamen concipienda est quaedam materiae\protect\index{Sachverzeichnis}{materia}
ad motum resistentia,\protect\index{Sachverzeichnis}{resistentia materiae ad motum}
quae etiam ab initio assumta restitantiam\protect\index{Sachverzeichnis}{restitantia} daret,
adeoque vim efficeret,
perinde ac si corpora 
\edtext{aequali vi restitiva\protect\index{Sachverzeichnis}{vis restitiva} occurrerent;}{%
\lemma{aequali}\Bfootnote{%
\textit{(1)}~vi respectiva concurrerent
\textit{(2)}~\textbar~concur- \textit{streicht Hrsg.}~\textbar\
\textit{(3)}~vi restitiva occurrerent;%
~\textit{L}}}
%
interim distincte explicando causas,\protect\index{Sachverzeichnis}{causa explicanda}%
\edlabel{LH_35_10_06_037r_blstft_udhakv}
remque adeo reducendo ad mutationes insensibiles,\protect\index{Sachverzeichnis}{mutatio insensibilis}
ut evitetur mutatio per saltum,\protect\index{Sachverzeichnis}{mutatio per saltum}
via quam proposui optima distinctissimaque est\lbrack;\rbrack\
etsi enim materia nostri gradus\protect\index{Sachverzeichnis}{materia nostri gradus}
sit diversae densitatis\protect\index{Sachverzeichnis}{densitas materiae}
quoad alia fluida,
\edtext{attamen quoad}{%
\lemma{attamen}\Bfootnote{%
\textit{(1)}~quod
\textit{(2)}~quoad%
~\textit{L}}}
%
fluidum suum elasticificans\protect\index{Sachverzeichnis}{fluidum elasticificans} ubique uniformiter disgregata,
seu in partes aequales divisa concipi potest.
Cum fluida nobis tenuia,\protect\index{Sachverzeichnis}{fluidum tenue}
huic tenuissimo rursus sint crassa,\protect\index{Sachverzeichnis}{fluidum tenue}
et pertusa,
ut ipsius respectu caetera uniforme quid esse videantur\lbrack,\rbrack\
ut si tam globi\protect\index{Sachverzeichnis}{globus minimus} majores
quam minores minimis 
\edtext{aequalibus foraminibus\protect\index{Sachverzeichnis}{foramen} pertusi,}{%
\lemma{aequalibus}\Bfootnote{%
\textit{(1)}~pertusi p
\textit{(2)}~foraminibus pertusi,%
~\textit{L}}}
%
vel
\edtext{potius in globos sui generis minimos\protect\index{Sachverzeichnis}{globus minimus} aequales}{%
\lemma{potius}\Bfootnote{%
\textit{(1)}~in aequales
\textit{(2)}~in globos \lbrack...\rbrack\ minimos aequales%
~\textit{L}}}
%
divisi intelligantur.
\pend%
%
%
%  \newpage
  \vspace{1.5em}%
  \centerline{\includegraphics[width=0.26\textwidth]{gesamttex/edit_VIII,3/images/LH_35_10_06_037_d.pdf}}%
  \vspace{0.0em}%
  \centerline{\lbrack\textit{Fig.~1}\rbrack}% \hspace*{16mm}
  \label{LH_35_10_06_037r_Fig.1}%
  \vspace{1.5em}%
%
%
\pstart%
Sunt autem in Systematibus\protect\index{Sachverzeichnis}{systema} concipiendi
velut gradus.\protect\index{Sachverzeichnis}{gradus}
Sic hi globuli\protect\index{Sachverzeichnis}{globulus} minimi nostri Systematis
fortasse sint instar planetarum \protect\index{Sachverzeichnis}{planeta}
\edtext{Systematis inferioris.\protect\index{Sachverzeichnis}{systema inferius}}{%
\lemma{Systematis}\Bfootnote{%
\textit{(1)}~major
\textit{(2)}~inferioris.%
~\textit{L}}}
%
\pend
\count\Bfootins=1200%
\count\Afootins=1200%
\count\Cfootins=1200%
%
%
% Ende des Stückes auf Bl. 37 r.